% Internal structural theorem. This does not prove general core unimodality.
\section{Cumulative reverse deficits give individual window ratios}
For an increasing arm $T$ of weight $c$, write
\[
 E_T=1-z+(2z-1)F_T,\quad R_T=z^{c+1}E_T(1/z),\quad
 C_T(n)=\sum_{j=0}^n[z^j]R_T.
\]
All cumulative sequences are zero at negative indices and constant
after their finite input ends. The empty arm has $R_\varnothing=1$.

\begin{theorem}\label{thm:deficit-cumulative-lc}
For every increasing arm, $C_T$ is log-concave on all integer indices.
\end{theorem}
\begin{proof}
Induct using $R_{Tt}=[t]R_T+z^c$. The empty case is constant one
on the nonnegative integers. For a nonempty prefix the positive
support of $R_T$ is $[0,d]$, where $d=c+1-\kappa\le c-1$;
the increasing condition gives $\kappa\le t+1$. Its mass $m$ equals
$F_T(1)\ge c+1$, since every rank from zero to $c$ is nonempty.

Set $w=[t]R_T$ and $v_n=\sum_{j=0}^n w_j=[t]C_T(n)$.
By induction and log-concave convolution, $v$ is log-concave.
The new cumulative sequence is $h_n=v_n+\mathbf1_{n\ge c}$.
We check the joining edge and subsequent centers.

If $c\ge t$, then
\[
 w_{c-1}-w_c=R_T(c-t)\ge1.
\]
Indeed $0\le c-t\le d$. The known unimodality of $R_T$
and of a uniform window makes $w$ nonincreasing from $c-1$
onward. Thus $v$ is concave there. Writing
$V=v_{c-2}$, $x=w_{c-1}$ and $y=w_c$, the determinant at
the only left joining center is
\[
 h_{c-1}^2-h_{c-2}h_c=V(x-y-1)+x^2\ge0.
\]

If $1\le c<t$, then $w_n=m$ on $[c-1,t-1]$ and is
nonincreasing afterward. Again $v$ is concave from $c-1$.
At the joining center its determinant is
\[
 m^2-v_{c-2}\ge m^2-(c-1)m\ge0.
\]
In both cases the determinant at $c$ is the old determinant
plus $2v_c+1-v_{c-1}\ge0$. At every center $n\ge c+1$ it is
the old determinant plus $2v_n-v_{n-1}-v_{n+1}\ge0$.
Earlier centers are unchanged. Finally, for the empty prefix
$c=0$, the new cumulative sequence is $2,3,\ldots,t+1$ followed
by a constant tail, which is log-concave. This closes induction.
\end{proof}

\begin{lemma}[A cumulative premise implies actual prefix-to-tail ratios]
Let $p$ be a finite nonnegative sequence, $p_0>0$, and suppose
its cumulative sequence $C_n$ is log-concave. For every $t\ge2$,
the actual coefficients $w=[t]p$ satisfy
\[
 w_{k+1}w_j\ge w_kw_{j+1}\qquad(0\le k\le t-2,\ j\ge k).
\]
\end{lemma}
\begin{proof}
Fix $k$ and put $r=C_{k+1}/C_k$. Since $w_k=C_k$ and
$w_{k+1}=C_{k+1}$, it suffices to show $w_{j+1}\le rw_j$.
For $k\le j\le k+t$, cumulative log-concavity gives
$C_{j+1}\le rC_j$ and $C_{j+1-t}\ge rC_{j-t}$;
zero extension covers negative lower indices. Subtraction proves
the claim. For $i\ge k+1$, the same premise gives
\[
 p_{i+1}C_{i-1}\le p_iC_i,\qquad
 p_{i+1}\le (C_i/C_{i-1})p_i\le rp_i.
\]
Summing these inequalities over the window proves the remaining
case $j\ge k+t$. All denominators used are positive.
\end{proof}

\begin{corollary}\label{cor:deficit-window-ratios}
For every increasing arm $T$ and every $t\ge2$, $[t]R_T$
has the prefix-to-tail coefficient comparison through $t-2$.
The window need not exceed the largest part of $T$.
\end{corollary}

\begin{proposition}[The heavier one-part terminal case at weight two]
Let $A=F_{1^\alpha t}$ and let $B$ be any nonrectangular
increasing-arm polynomial. Suppose $t\ge2$, $\alpha\ge1$,
and $a=\alpha+t\ge b=\deg B$. For $q=z^{a+b}(AB)(1/z)$,
the sequence $2^j\max(q_j-q_{j-1},0)$ is unimodal.
\end{proposition}
\begin{proof}
The identity $A=1+z[t][\alpha+1]$ permits exchange of the
two interval factors. We may therefore assume $\alpha\ge t-1$.
The exchanged capacity is smaller if an exchange was needed.
Put $q_B=z^bB(1/z)$ and $V=[t]q_B$. Then
\[
 q=[\alpha+1]V+z^a q_B.
\]
Its first mode $s$ is strictly less than $a$, by ordinary-product
first-half growth and the extra midpoint rank for equal nonrectangular
arms. Through $s$, its positive differences are
$d_j=V_j-V_{j-\alpha-1}$. They are $V_j$ for $j\le\alpha$.
The sequence $2^jV_j$ is nondecreasing: in forward order the
coefficients of $[t]B$ grow by a factor at most two. For ranks
below $t$ use $f_n\le\sum_{i<n}f_i$; afterward use
$f_n\le2f_{n-1}$ termwise.

For $\alpha+1\le j\le s$, set $k=j-\alpha-1\le t-2$.
Also $s\le\deg V=b+t-1$, so $k<b$ by normalization.
Writing $W=[t]R_B$, the exact identity is
\[
 (2-z)V=W+z^b-z^{b+t},\qquad
 2d_j-d_{j-1}=W_j-W_k+\mathbf1_{j=b}.
\]
The negative terminal monomial is outside $j\le s$ and neither
monomial affects the head index $k<b$. By the preceding corollary,
$W_{k+\alpha+1}/W_k$ is nonincreasing wherever its denominator
is positive. Thus $W_j-W_k$ has no negative-to-positive change.
If such a difference was already negative before $j=b$, it is
at most $-1$ at $b$, by integrality and the ordered ratios; the
extra unit still cannot make it positive. Zero endpoints are
handled directly. The actual weighted differences have the same
single-crossing property. Joining to their nondecreasing initial
segment and then to terminal zeros proves the proposition.
\end{proof}

These results settle a terminal case of the proposed general
$2^j$-difference induction. Closure under longer uniform windows
still requires a separate argument; the unrestricted conjecture
is not proved. Cumulative log-concavity is used here only to prove
the displayed inequalities for individual window coefficients.
